\section{二维双调和函数}
\label{sec:phys-cont-biharmonic-primer}

梁上的四阶导数换到平面时，最简单的对应物是拉普拉斯算子的平方。设 $\Omega\subset\mathbb R^2$ 为开集，$u\in C^4(\Omega)$。若 $\Delta^2u:=\Delta(\Delta u)=0$，便称 $u$ 为双调和函数；满足 $\Delta u=0$ 的函数称为调和函数。每个调和函数当然双调和，反过来却不成立：$u(x,y)=x^2+y^2$ 满足 $\Delta u=4$、$\Delta^2u=0$。因此也不能直接套用调和函数的最大值原理。例如单位圆盘上的 $1-r^2$ 在边界为零、内部为正，同时仍是双调和函数。

给定载荷 $f$ 时，方程 $D\Delta^2u=f$ 只有内部条件；在边界上还须给出两项相互独立的数据。夹持边界给 $u=\partial_nu=0$。数学上常用的 Navier 边界给 $u=\Delta u=0$；对各向同性薄板的直边段，由于边界上 $u=0$ 导致切向二阶导数为零，它与零弯矩的简支条件一致，但在弯曲边界上不能不加检验地等同。自由边界的两项条件是弯矩和有效剪力为零，稍后的变分公式会把它们准确写出。只给 $u=0$ 并不足以构成四阶边值问题；上面的 $1-r^2$ 已提供一个非零齐次解。

先看圆形几何为何总出现对数。若 $u=u(r)$ 不依赖极角，把 $\partial_xr=x/r$、$\partial_yr=y/r$ 代入链式法则，可得
\begin{equation}
\Delta u=u''(r)+\frac1r u'(r)=\frac1r(ru')',\qquad r>0.
\label{eq:phys-cont-radial-laplacian}
\end{equation}
令 $g=\Delta u$。$\Delta^2u=0$ 意味着 $(rg')'=0$，所以 $g=A\log r+B$。再解 $(ru')'=r(A\log r+B)$，积分两次得到
\begin{equation}
u(r)=C_1+C_2r^2+C_3\log r+C_4r^2\log r,\qquad 0<r<R.
\label{eq:phys-cont-radial-biharmonic-basis}
\end{equation}
四个系数对应四阶方程的四维局部解空间。若区域包含原点且要求 $u\in C^2$，必须令 $C_3=C_4=0$：$\log r$ 本身发散；$r^2\log r$ 虽趋于零，其二阶导数仍含 $\log r$。在圆环里没有原点，后两项则完全合法。后面圆环薄板的对数解不是凑出的试探函数，而是解空间不可少的成员。

对非轴对称问题也可分离变量。由 $x=r\cos\theta$、$y=r\sin\theta$ 的链式法则先得
\[
\partial_x=\cos\theta\,\partial_r-\frac{\sin\theta}{r}\partial_\theta,
\qquad
\partial_y=\sin\theta\,\partial_r+\frac{\cos\theta}{r}\partial_\theta.
\]
将两式各作用两次并相加，可核算 $\Delta=\partial_r^2+r^{-1}\partial_r+r^{-2}\partial_\theta^2$。因此
\[
\Delta(r^p\ee^{\ii m\theta})=(p^2-m^2)r^{p-2}\ee^{\ii m\theta},
\]
其中 $m=0,1,2,\ldots$。对 $m\ge1$，圆心正则的两支双调和模态是 $r^m\ee^{\ii m\theta}$ 与 $r^{m+2}\ee^{\ii m\theta}$；第二支第一次作用 $\Delta$ 后成为调和的第一支，第二次便为零。对 $m=0$，正则两支是 $1,r^2$。更一般的函数也有同样的“两个调和部分”，而且不必假定它能先展开成 Fourier 级数。

\begin{theorem}[双调和函数的 Almansi 分解]
设 $\Omega\subset\mathbb R^2$ 对原点星形，即 $x\in\Omega$ 时整条线段 $\{tx:0\le t\le1\}$ 都在 $\Omega$ 内。若 $u\in C^4(\Omega)$ 且 $\Delta^2u=0$，则存在唯一一对调和函数 $h_0,h_1$，使
\begin{equation}
u(x)=h_0(x)+|x|^2h_1(x),\qquad \Delta h_0=\Delta h_1=0.
\label{eq:phys-cont-almansi-polynomial}
\end{equation}
\end{theorem}

\begin{proof}
记 $g=\Delta u$。由于 $\Delta g=0$，定义
\[
h_1(x)=\frac14\int_0^1g(tx)\dd t.
\]
对 $x$ 求 Laplace 时，每次导数都带一个 $t$，所以 $\Delta h_1=\frac14\int_0^1t^2(\Delta g)(tx)\dd t=0$。令 $E=x\partial_x+y\partial_y$ 为径向缩放导数。乘积法则直接给出 $\Delta(|x|^2h)=4h+4Eh+|x|^2\Delta h$。另一方面，沿线段对 $t$ 作一次分部积分，得到
\[
(I+E)h_1(x)
=\frac14\int_0^1\frac{\dd}{\dd t}[t g(tx)]\dd t
=\frac14g(x).
\]
故 $\Delta(|x|^2h_1)=g=\Delta u$，从而 $h_0:=u-|x|^2h_1$ 调和。若另有一组这样的分解，对两式取 Laplace 后，其 $h_1$ 的差满足 $(I+E)h_1=0$；沿每条射线积分 $\frac{\dd}{\dd t}[t h_1(tx)]=0$，再用原点处的连续性，得差为零，$h_0$ 的差也为零。
\end{proof}

即使不使用完整分解，也能求出圆周平均。设闭圆盘 $\overline{B_R(0)}\subset\Omega$，记 $A(r)=(2\pi)^{-1}\int_0^{2\pi}u(r,\theta)\dd\theta$。把极坐标拉普拉斯中的 $r^{-2}\partial_\theta^2u$ 沿一周积分，因周期性它为零，故 $A$ 仍满足径向双调和方程。$u$ 在圆心光滑，排除\eqnrefs{eq:phys-cont-radial-biharmonic-basis}中的对数项。再由 $A(0)=u(0)$ 和 $\Delta A(0)=\Delta u(0)$ 确定系数，得到
\begin{equation}
\frac1{2\pi}\int_0^{2\pi}u(r,\theta)\dd\theta
=u(0)+\frac{r^2}{4}\Delta u(0),\qquad 0<r<R.
\label{eq:phys-cont-biharmonic-mean-value}
\end{equation}
与调和函数的平均值性质相比，多出的 $r^2\Delta u(0)/4$ 正是两者的区别。

点载荷带来一种须在分布意义下理解的双调和函数。在整个平面定义
\begin{equation}
\Phi(x)=\frac{|x|^2\log|x|}{8\pi},\qquad x\ne0.
\label{eq:phys-cont-biharmonic-fundamental}
\end{equation}
由\eqnrefs{eq:phys-cont-radial-laplacian}算得 $\Delta(r^2\log r)=4\log r+4$。记 $\delta_0$ 为原点的点质量，即对任意光滑紧支撑试验函数 $\varphi$，$\int\delta_0(x)\varphi(x)\dd^2x=\varphi(0)$。另一方面，$\Delta\log r=2\pi\delta_0$：圆心外它为零，而半径 $\varepsilon$ 小圆上的通量 $\int\partial_r\log r\,\dd s=\int_0^{2\pi}(1/\varepsilon)\varepsilon\dd\theta=2\pi$。由于 $r^2\log r$ 及其一阶导数在小圆上的边界贡献趋零，第一次取 Laplace 不另产生点质量。于是
\begin{equation}
\Delta^2\Phi=\delta_0
\label{eq:phys-cont-biharmonic-fundamental-delta}
\end{equation}
作为分布成立。在有界区域，$\Phi(x-y)$ 通常不满足夹持或简支边界条件，不能直接称为该边值问题的 Green 函数；必须再加一个关于 $x$ 双调和的修正项，使两项边界条件同时成立。

夹持边值问题的唯一性也可直接证明。若 $\Delta^2w=0$ 且边界上 $w=\partial_nw=0$，两次分部积分消去边界项，得到
\[
0=\int_\Omega w\Delta^2w\dd^2x=\int_\Omega(\Delta w)^2\dd^2x.
\]
所以 $\Delta w=0$；再由 $w=0$ 的边界数据与调和函数的能量恒等式，得 $\int_\Omega|\nabla w|^2=0$，从而 $w=0$。这证明的是唯一性；存在性还需要强制变分理论，而不是由分部积分自动得到。

最后用一个能手算也能数值复核的载荷收束本节。单位正方形 $\Omega=(0,1)^2$ 上取 Navier 边界和 $f(x,y)=\sin(\pi x)\sin(\pi y)$。因为 $\Delta f=-2\pi^2f$，故 $\Delta^2f=4\pi^4f$；于是
\begin{equation}
\Delta^2u=f,\qquad u|_{\partial\Omega}=\Delta u|_{\partial\Omega}=0
\quad\Longrightarrow\quad
u(x,y)=\frac{\sin(\pi x)\sin(\pi y)}{4\pi^4}.
\label{eq:phys-cont-square-biharmonic-exact}
\end{equation}
解的唯一性可令 $v=\Delta u$，对 $\Delta v=f$、$\Delta u=v$ 两个零 Dirichlet 问题依次应用二阶能量法。\figref{fig:phys-cont-biharmonic-square}把精确解与五点差分 Laplace 算子连续作用两次所得的离散解比较。连续方程的两次微分在数值上变成两次求解稀疏线性系统，而网格误差随步长平方下降。

\begin{figure}[htbp]
\centering
\includegraphics[width=.94\linewidth]{biharmonic_square.pdf}
\caption{单位正方形上 $\Delta^2u=\sin(\pi x)\sin(\pi y)$ 的 Navier 解。上图为数值挠度，下图为相对精确式\eqnrefs{eq:phys-cont-square-biharmonic-exact}的最大误差；采用内点五点差分与两次零边界 Poisson 求解。}
\label{fig:phys-cont-biharmonic-square}
\end{figure}
